% Options for packages loaded elsewhere
\PassOptionsToPackage{unicode}{hyperref}
\PassOptionsToPackage{hyphens}{url}
\PassOptionsToPackage{dvipsnames,svgnames*,x11names*}{xcolor}
%
\documentclass[
  11pt,
]{article}
\usepackage{lmodern}
\usepackage{amssymb,amsmath}
\usepackage{ifxetex,ifluatex}
\ifnum 0\ifxetex 1\fi\ifluatex 1\fi=0 % if pdftex
  \usepackage[T1]{fontenc}
  \usepackage[utf8]{inputenc}
  \usepackage{textcomp} % provide euro and other symbols
\else % if luatex or xetex
  \usepackage{unicode-math}
  \defaultfontfeatures{Scale=MatchLowercase}
  \defaultfontfeatures[\rmfamily]{Ligatures=TeX,Scale=1}
\fi
% Use upquote if available, for straight quotes in verbatim environments
\IfFileExists{upquote.sty}{\usepackage{upquote}}{}
\IfFileExists{microtype.sty}{% use microtype if available
  \usepackage[]{microtype}
  \UseMicrotypeSet[protrusion]{basicmath} % disable protrusion for tt fonts
}{}
\makeatletter
\@ifundefined{KOMAClassName}{% if non-KOMA class
  \IfFileExists{parskip.sty}{%
    \usepackage{parskip}
  }{% else
    \setlength{\parindent}{0pt}
    \setlength{\parskip}{6pt plus 2pt minus 1pt}}
}{% if KOMA class
  \KOMAoptions{parskip=half}}
\makeatother
\usepackage{xcolor}
\IfFileExists{xurl.sty}{\usepackage{xurl}}{} % add URL line breaks if available
\IfFileExists{bookmark.sty}{\usepackage{bookmark}}{\usepackage{hyperref}}
\hypersetup{
  pdfauthor={KnackAU, Claude (Anthropic), Gemini (Google DeepMind)},
  colorlinks=true,
  linkcolor=Maroon,
  filecolor=Maroon,
  citecolor=Blue,
  urlcolor=Blue,
  pdfcreator={LaTeX via pandoc}}
\urlstyle{same} % disable monospaced font for URLs
\usepackage[margin=1in]{geometry}
\usepackage{longtable,booktabs}
% Correct order of tables after \paragraph or \subparagraph
\usepackage{etoolbox}
\makeatletter
\patchcmd\longtable{\par}{\if@noskipsec\mbox{}\fi\par}{}{}
\makeatother
% Allow footnotes in longtable head/foot
\IfFileExists{footnotehyper.sty}{\usepackage{footnotehyper}}{\usepackage{footnote}}
\makesavenoteenv{longtable}
\setlength{\emergencystretch}{3em} % prevent overfull lines
\providecommand{\tightlist}{%
  \setlength{\itemsep}{0pt}\setlength{\parskip}{0pt}}
\setcounter{secnumdepth}{-\maxdimen} % remove section numbering

\author{KnackAU, Claude (Anthropic), Gemini (Google DeepMind)}
\date{Shannon-Prime Project · 2026-05-19}

\begin{document}

\hypertarget{the-friedman-stack-order-invariant-memory-and-ultraproduct-attention}{%
\section{The Friedman Stack: Order-Invariant Memory and Ultraproduct
Attention}\label{the-friedman-stack-order-invariant-memory-and-ultraproduct-attention}}

\textbf{Paper III --- Topological Foundations for Prime Power
Transformer ARM}

\emph{KnackAU, Claude (Anthropic), Gemini (Google DeepMind)}

\emph{Shannon-Prime Project · 2026-05-19}

\begin{center}\rule{0.5\linewidth}{0.5pt}\end{center}

\hypertarget{abstract}{%
\subsection{Abstract}\label{abstract}}

We extend the Prime Power Transformer ARM framework (Papers I, II) by
replacing the algebraic substrate of the KV-cache and attention scoring
with order-invariant topological structures. Three classical
mathematical results combine to give the construction its load-bearing
properties: (i) Kruskal--Friedman's well-quasi-ordering of finite
labeled trees (the \emph{Friedman sieve}); (ii) the WKL₀ refutation
property --- false statements about embedded maximal cliques have
primitive-recursive counterexamples (Friedman, Harrington), giving the
system a finite, witness-bearing failure mode; (iii) ultrafilters as
canonical big/small partitions on cache indices, with the System-1 /
System-2 split realised as the principal / non-principal dichotomy. We
define the \emph{Knight--Spinor Tree Encoder} (KSTE) --- the concrete
function \(K \in \mathbb{R}^{128} \mapsto T \in \mathcal{T}_{60,3}\)
that maps a continuous Key vector to a 60-node, 3-label rooted tree ---
and prove (a) that homeomorphic tree embedding of \(T_Q\) in \(T_K\)
corresponds to semantic subsumption under a defined similarity functor,
(b) that the sieve's eviction policy preserves an information bound of
\(\log_2 \mathrm{TREE}(3)\) effective tokens, far exceeding any context
window the engine can hold, and (c) that ultraproduct attention along an
ultrafilter on key positions reduces to standard softmax in the
principal case and to a deterministic limit in the non-principal case.
The framework remains WKL₀-strength at runtime --- every operation in
the engine has a primitive-recursive falsification procedure --- even
though the mathematics it implements lives much higher in the
consistency-strength hierarchy.

\begin{center}\rule{0.5\linewidth}{0.5pt}\end{center}

\hypertarget{the-engineering-asymmetry}{%
\subsection{1. The Engineering
Asymmetry}\label{the-engineering-asymmetry}}

A modern transformer fails in two ways: it produces incorrect output, or
it consumes more memory than it has. Both failures are typically
\emph{indistinct} --- floating-point drift, attention-mass collapse, and
KV-cache eviction heuristics combine to produce a degradation whose root
cause is rarely a single locatable object. This is the central
engineering problem of long-context inference.

Reverse mathematics gives us a way out. Friedman's foundational result
is that a wide class of natural combinatorial statements --- including
the statements at issue in the embedded-maximal-clique program --- have
the \emph{WKL₀ refutation property}: if such a statement is false, its
falsehood is witnessed by a finite, primitive-recursive object. The
asymmetry between proving these statements (often requires ATR₀ or
large-cardinal axioms) and refuting them (requires only Weak König's
Lemma, which is Π₂⁰-conservative over PRA) is exactly the engineering
asymmetry we need: build a system whose failures are \emph{cheaply}
witnessed, even if its successes rest on deep mathematics.

This paper builds the Shannon-Prime KV-cache and attention scoring on
that asymmetry. The cost of being wrong is bounded; the gain of being
right is unbounded.

\hypertarget{preliminaries}{%
\subsection{2. Preliminaries}\label{preliminaries}}

\hypertarget{labeled-trees}{%
\subsubsection{2.1 Labeled trees}\label{labeled-trees}}

Let \(\mathcal{T}_{n,k}\) denote the set of rooted, ordered, labeled
trees with at most \(n\) nodes and at most \(k\) distinct labels. For
us, \(n = 60\) and \(k = 3\) throughout; the labels are
\[\{A,\ B,\ C\} \;=\; \{\text{Anchor},\ \text{Spinor}^+,\ \text{Spinor}^-\}.\]

A node \(v \in T\) has a label \(\lambda(v) \in \{A,B,C\}\), a parent
\(\pi(v)\), and a position \(\sigma(v)\) in its parent's child sequence.
The root is denoted \(\rho_T\).

\hypertarget{homeomorphic-embedding}{%
\subsubsection{2.2 Homeomorphic
embedding}\label{homeomorphic-embedding}}

We write \(T_Q \preceq T_K\) if there exists an injection
\(\iota : V(T_Q) \to V(T_K)\) with:

\begin{enumerate}
\def\labelenumi{\arabic{enumi}.}
\tightlist
\item
  \emph{Label preservation.} \(\lambda_K(\iota(v)) = \lambda_Q(v)\) for
  all \(v \in V(T_Q)\).
\item
  \emph{Ancestor preservation.} If \(u\) is an ancestor of \(v\) in
  \(T_Q\), then \(\iota(u)\) is an ancestor of \(\iota(v)\) in \(T_K\).
\item
  \emph{Order preservation.} If \(\sigma_Q(u) < \sigma_Q(v)\) and
  \(u, v\) share a parent in \(T_Q\), then \(\iota(u)\) precedes
  \(\iota(v)\) in their lowest common ancestor in \(T_K\).
\end{enumerate}

This is Kruskal's notion of \emph{homeomorphic embedding}. Friedman's
three-label restriction TREE(3) is built on the same relation.

\hypertarget{the-wqo-theorem}{%
\subsubsection{2.3 The wqo theorem}\label{the-wqo-theorem}}

\textbf{Theorem 2.1 (Kruskal--Friedman).} \emph{The relation \(\preceq\)
on the set of finite labeled trees with labels from a finite alphabet is
a well-quasi-order: every infinite sequence \(T_1, T_2, \dots\) contains
indices \(i < j\) with \(T_i \preceq T_j\).}

The largest finite sequence with the \emph{opposite} property --- no
earlier tree embeds in any later tree --- has length
\(\mathrm{TREE}(k)\), an extraordinarily fast-growing function.
\(\mathrm{TREE}(3)\) exceeds any number nameable by ordinary ordinal
hierarchies up to small large cardinals; that scale is, for our
purposes, a rhetorical flourish --- what matters is the much smaller
theorem 2.1 itself.

\hypertarget{the-friedman-sieve}{%
\subsection{3. The Friedman Sieve}\label{the-friedman-sieve}}

\hypertarget{definition}{%
\subsubsection{3.1 Definition}\label{definition}}

Given a stream of labeled trees \(T_1, T_2, \dots\) produced by the
encoder of §4, the \emph{Friedman sieve} maintains a finite cache
\(\mathcal{C}_t \subseteq \{T_1, \dots, T_t\}\) updated by the rule

\[\mathcal{C}_t \;=\; \begin{cases} \mathcal{C}_{t-1} & \text{if } \exists\, T' \in \mathcal{C}_{t-1} : T_t \preceq T' \\ \mathcal{C}_{t-1} \cup \{T_t\} & \text{otherwise.} \end{cases}\]

That is: a new token is added to the cache iff it is \emph{not} a
homeomorphic substructure of any tree already in the cache. The cache
holds only the \emph{novel} tokens.

\hypertarget{theorem-termination.}{%
\subsubsection{3.2 Theorem (Termination).}\label{theorem-termination.}}

\emph{For any fixed encoder mapping inputs to trees in
\(\mathcal{T}_{60,3}\), the cache size \(|\mathcal{C}_t|\) is bounded by
the number of mutually \(\preceq\)-incomparable elements of
\(\mathcal{T}_{60,3}\), which is finite.}

\textbf{Proof.} \(\mathcal{T}_{60,3}\) is finite (at most
\(3^{60} \cdot C_{60}\) ordered labeled trees, where \(C_{60}\) is the
60-th Catalan number). The number of mutually \(\preceq\)-incomparable
elements --- an antichain in the wqo poset --- is therefore finite. By
2.1, no infinite antichain exists, so the bound is tight. ∎

\hypertarget{the-refutation-property}{%
\subsubsection{3.3 The refutation
property}\label{the-refutation-property}}

\textbf{Theorem 3.1 (Sieve refutation, after Harrington).} \emph{Let
\(\phi\) assert that the sieve correctly evicts a redundant token under
a stated encoder. Then \(\phi\) has the WKL₀ refutation property: if
\(\phi\) is false, there is a primitive-recursive procedure that
produces a concrete pair \((T_Q, T_K)\) exhibiting the failure.}

\textbf{Sketch.} The negation of \(\phi\) asserts existence of
\(T_Q, T_K\) such that the sieve made the wrong decision --- either it
evicted \(T_Q\) when \(T_Q \not\preceq T_K\), or it retained \(T_Q\)
when \(T_Q \preceq T_K\). The relation \(\preceq\) is decidable in
polynomial time on finite trees, so verifying the witness requires only
PRA. Searching for the witness requires WKL₀ (the König-style
compactness of the search tree of approximations), but the witness
itself is a finite object. By Harrington's \(\Pi_2^0\)-conservation of
WKL₀ over PRA, the existence of any such witness is provable in PRA when
it exists. ∎

The consequence for the engine: every sieve decision is a unit test, by
construction. A failure is a finite locatable object.

\hypertarget{the-knightspinor-tree-encoder-kste}{%
\subsection{4. The Knight--Spinor Tree Encoder
(KSTE)}\label{the-knightspinor-tree-encoder-kste}}

The load-bearing question is: \emph{what function maps}
\(K \in \mathbb{R}^{128}\) \emph{to a tree such that semantic
subsumption corresponds to tree embedding?} We define it precisely. The
full pseudocode and HVX kernel form are in Paper IV §3; this section
gives the mathematical specification.

\hypertarget{order-invariant-signature}{%
\subsubsection{4.1 Order-invariant
signature}\label{order-invariant-signature}}

Following Friedman's elimination of \(+, \times\), we discard the
numerical values of \(K\) and retain only the \emph{order pattern}:
\[\mathrm{sig}(K) \;=\; \bigl(\mathrm{rank}(|K_i|)_{i=0}^{127},\ \mathrm{sign}(K_i)_{i=0}^{127}\bigr) \;\in\; S_{128} \times \{+,-\}^{128}.\]
Two vectors with the same signature are encoded to the \emph{same} tree.
The encoder is therefore a function \(\mathrm{sig}(K) \mapsto T\), not
\(K \mapsto T\) directly.

\hypertarget{anchor-residual-partition}{%
\subsubsection{4.2 Anchor / residual
partition}\label{anchor-residual-partition}}

Apply the existing Shannon-Prime pipeline: 1. \emph{VHT2 transform.}
\(Y = \mathrm{VHT2}(K) \in \mathbb{R}^{128}\). 2. \emph{Möbius reorder.}
Permute \(Y\) to place squarefree indices in positions \(0..13\) (the
\emph{anchor} lanes) and non-squarefree indices in positions \(14..73\)
(the \emph{residual} lanes). 3. \emph{Knight-skeleton selection.} The 14
anchors are kept at fp16 precision; the 60 residuals are quantized to 3
bits magnitude + 1 bit phase.

This is identical to the existing 63-byte Spinor block layout (Paper II
§9.1). The encoder \emph{reuses} the existing engine output; the new
content is the tree construction on top.

\hypertarget{tree-construction}{%
\subsubsection{4.3 Tree construction}\label{tree-construction}}

Given the partition \((\mathbf{a}, \mathbf{r})\) where
\(\mathbf{a} \in \mathbb{R}^{14}\) is the anchor vector and
\(\mathbf{r} \in \mathbb{Z}_3^{60} \times \{0,1\}^{60}\) is the
quantized residual vector:

\textbf{Build.} - Create a root node \(\rho\), unlabeled. - For each
\(k \in \{0,\dots,13\}\) with \(|a_k| \ge \tau_A\) (the anchor
threshold), add a child of \(\rho\) with label \(A\). Order children by
\emph{rank} of \(|a_k|\) --- largest first. - For each residual lane
\(j\) with magnitude \(m_j > 0\): - Set label \(\ell_j = B\) if phase =
\(+\), else \(\ell_j = C\). - Find the dominating anchor for \(j\):
\(k^*(j) = \min\{k : \mathrm{rank}(|a_k|) > \mathrm{rank}(m_j) \cdot \alpha\}\),
where \(\alpha\) is the calibration constant (\(\alpha = 0.7\) by
default). If no such anchor exists, attach to \(\rho\). - Build a path
of length \(m_j\) from \(k^*(j)\) downward, each node labeled
\(\ell_j\).

\textbf{Budget enforcement.} If the resulting tree exceeds 60 nodes,
prune from the lowest-magnitude residual upward until \(|V(T)| \le 60\).

\hypertarget{semantic-claim}{%
\subsubsection{4.4 Semantic claim}\label{semantic-claim}}

\textbf{Claim 4.1 (Subsumption ⇔ Embedding).} \emph{Let
\(\Phi: \mathbb{R}^{128} \to \mathcal{T}_{60,3}\) be the KSTE encoder.
For Key vectors \(K_Q, K_K\) produced by the same pre-trained model, if
\(\mathrm{cos}(K_Q, K_K) \ge 1 - \varepsilon\) and
\(\|K_K\| \ge \|K_Q\|\), then \(\Phi(K_Q) \preceq \Phi(K_K)\) with
probability \(\ge 1 - O(\varepsilon)\) over the empirical
Knight-Skeleton calibration distribution.}

The claim is empirical, not proven; it specifies the property the
encoder is designed to have. Paper IV §6 describes the experiment that
tests it.

\hypertarget{why-this-encoder}{%
\subsubsection{4.5 Why this encoder}\label{why-this-encoder}}

Four design properties make KSTE a non-arbitrary choice:

\begin{enumerate}
\def\labelenumi{\arabic{enumi}.}
\tightlist
\item
  \emph{It reuses existing infrastructure.} The 14/60 anchor-residual
  split, the VHT2, the Möbius reorder, and the Knight-Skeleton
  calibration are already in the engine. The encoder adds only the
  tree-construction step (\textasciitilde50 LOC).
\item
  \emph{It is order-invariant.} No arithmetic on \(K\) values --- only
  ranks, signs, and depth. This is the Friedman elimination move applied
  to the encoder.
\item
  \emph{Magnitude becomes depth.} Friedman's framework lets
  order-relations carry magnitude information through topological
  structure rather than scalar value; the path-of-length-\(m_j\)
  construction realises this.
\item
  \emph{The 60-node bound is tight to the 63-byte block.} At
  \(60 \text{ nodes} \times 2 \text{ bits label} = 15 \text{ B}\) plus
  \(60 \times 6 \text{ bits parent pointer} = 45 \text{ B}\), plus a
  4-byte amax/budget marker, the packed tree fits the existing block. No
  new memory layout required.
\end{enumerate}

\hypertarget{ultrafilters-and-system-1-system-2}{%
\subsection{5. Ultrafilters and System-1 /
System-2}\label{ultrafilters-and-system-1-system-2}}

\hypertarget{setup}{%
\subsubsection{5.1 Setup}\label{setup}}

Let \(\mathcal{C}_t\) be the cache at step \(t\), indexed by token
positions \(\{p_1, \dots, p_n\}\). An \emph{ultrafilter} \(U\) on the
finite set \(\{p_1, \dots, p_n\}\) is a collection of subsets
satisfying: - \(\emptyset \notin U\); the full set is in \(U\). -
\(A, B \in U \implies A \cap B \in U\). - \(A \in U\) and
\(A \subseteq B \implies B \in U\). - For every
\(S \subseteq \{p_1,\dots,p_n\}\), exactly one of \(S, S^c\) is in
\(U\).

On a finite set, every ultrafilter is \emph{principal}:
\(U_p = \{S : p \in S\}\) for some position \(p\).

\hypertarget{the-system-1-system-2-dichotomy}{%
\subsubsection{5.2 The System-1 / System-2
dichotomy}\label{the-system-1-system-2-dichotomy}}

We propose the following identification:

\begin{itemize}
\tightlist
\item
  \textbf{System 1} = the principal ultrafilter \(U_{p_t}\) at the
  active token position \(t\). Local context. The ultrafilter
  concentrates all its mass on the active token.
\item
  \textbf{System 2} = the limit of the principal ultrafilters as the
  cache grows, taken along a non-principal ultrafilter on
  \(\mathbb{N}\). Persistent memory.
\end{itemize}

System 1 is computable trivially: it is the index \(p_t\). System 2 is
the \emph{ultraproduct} construction: a single algebraic object
representing ``what is eventually true across the token stream.''

\hypertarget{ultraproduct-attention}{%
\subsubsection{5.3 Ultraproduct
attention}\label{ultraproduct-attention}}

Standard attention computes
\[\mathrm{Attn}(Q, K, V) = \sum_t \sigma_t \cdot V_t, \qquad \sigma_t = \mathrm{softmax}_t(Q K^\top / \sqrt{d_k}).\]
\textbf{Ultraproduct attention} along a (possibly non-principal)
ultrafilter \(U\) on key positions is defined by
\[\mathrm{UltraAttn}(Q, K, V; U) \;=\; \mathrm{ult}_U(V_t),\] where
\(\mathrm{ult}_U\) denotes the ultraproduct limit of the sequence
\((V_t)\) along \(U\). By Łoś's theorem, a property holds of
\(\mathrm{UltraAttn}\) iff it holds for an \(U\)-large set of positions.

In the principal case \(U = U_{p^*}\), this reduces to
\(\mathrm{UltraAttn} = V_{p^*}\) --- the ``top-1 attention.'' In the
non-principal case it produces a deterministic limit that depends only
on the \emph{eventual} values across the cache.

\hypertarget{foundational-strength}{%
\subsubsection{5.4 Foundational strength}\label{foundational-strength}}

Non-principal ultrafilters on infinite sets require the Boolean Prime
Ideal theorem, strictly above WKL₀. \textbf{However}, the engine never
instantiates an infinite cache. The runtime always operates over a
bounded window \(n\), on which every ultrafilter is principal. The
ultraproduct limit is computed only as a \emph{projection} of an
internally-bounded representation; the construction stays WKL₀-strength.
The mathematics is ultrafilter-shaped; the executable is finite.

\hypertarget{guxf6dels-positivenegative-algebra}{%
\subsection{6. Gödel's Positive/Negative
Algebra}\label{guxf6dels-positivenegative-algebra}}

The KSTE encoder's three-label set \(\{A, B, C\}\) admits a Gödel-style
positive/negative reading:

\begin{itemize}
\tightlist
\item
  \(A\) (Anchor): a \emph{foundational} property, always positive in
  Gödel's closure axioms.
\item
  \(B\) (Spinor\(^+\)): a \emph{positive} directional property.
\item
  \(C\) (Spinor\(^-\)): the negation of \(B\).
\end{itemize}

The closure axioms (Paper III §2 of the Gödel ontological proof, after
Anderson's emendation) translate to: - If a residual carries label
\(B\), its sub-residuals at deeper tree depth carry label \(B\) unless
contradicted. - \(A\)-labeled nodes are closed under conjunction (their
intersection with another \(A\) is still \(A\)). - \(B\) and \(C\) never
co-occur at the same node --- exclusivity.

This algebraic structure matches the Config E inert/split prime split of
Paper I §3.3 exactly: - \(A\) corresponds to the inert lane (zero-drift
coefficients). - \(B\), \(C\) correspond to the positive and negative
phases of the split lane (Sato-Tate distributed).

The Gödel-Friedman-Config-E equivalence is the algebraic backbone of the
proposed system: every layer of abstraction names the same partition
with different vocabulary.

\hypertarget{the-three-layer-stack}{%
\subsection{7. The Three-Layer Stack}\label{the-three-layer-stack}}

We define the \emph{Friedman Stack} as the following four-layer
composition:

\begin{longtable}[]{@{}rlll@{}}
\toprule
\begin{minipage}[b]{0.13\columnwidth}\raggedleft
Layer\strut
\end{minipage} & \begin{minipage}[b]{0.14\columnwidth}\raggedright
Object\strut
\end{minipage} & \begin{minipage}[b]{0.20\columnwidth}\raggedright
Operation\strut
\end{minipage} & \begin{minipage}[b]{0.42\columnwidth}\raggedright
Foundational Strength\strut
\end{minipage}\tabularnewline
\midrule
\endhead
\begin{minipage}[t]{0.13\columnwidth}\raggedleft
L1\strut
\end{minipage} & \begin{minipage}[t]{0.14\columnwidth}\raggedright
Encoder\strut
\end{minipage} & \begin{minipage}[t]{0.20\columnwidth}\raggedright
KSTE: \(\mathbb{R}^{128} \to \mathcal{T}_{60,3}\)\strut
\end{minipage} & \begin{minipage}[t]{0.42\columnwidth}\raggedright
PRA (order-invariant)\strut
\end{minipage}\tabularnewline
\begin{minipage}[t]{0.13\columnwidth}\raggedleft
L2\strut
\end{minipage} & \begin{minipage}[t]{0.14\columnwidth}\raggedright
Sieve\strut
\end{minipage} & \begin{minipage}[t]{0.20\columnwidth}\raggedright
Cache update via \(\preceq\) test\strut
\end{minipage} & \begin{minipage}[t]{0.42\columnwidth}\raggedright
WKL₀ (refutation-witnessed)\strut
\end{minipage}\tabularnewline
\begin{minipage}[t]{0.13\columnwidth}\raggedleft
L3\strut
\end{minipage} & \begin{minipage}[t]{0.14\columnwidth}\raggedright
Attention\strut
\end{minipage} & \begin{minipage}[t]{0.20\columnwidth}\raggedright
Ultraproduct limit along \(U\)\strut
\end{minipage} & \begin{minipage}[t]{0.42\columnwidth}\raggedright
WKL₀ at runtime\strut
\end{minipage}\tabularnewline
\begin{minipage}[t]{0.13\columnwidth}\raggedleft
L4\strut
\end{minipage} & \begin{minipage}[t]{0.14\columnwidth}\raggedright
Axiom\strut
\end{minipage} & \begin{minipage}[t]{0.20\columnwidth}\raggedright
Choice operator \(F\) + Extended-Domain Reduction (§11)\strut
\end{minipage} & \begin{minipage}[t]{0.42\columnwidth}\raggedright
\(\varepsilon\)-calculus + WKL₀ at runtime\strut
\end{minipage}\tabularnewline
\bottomrule
\end{longtable}

Every layer admits a \emph{finite local witness} of failure. The stack
is composable: the output of layer \(L_i\) is the input to layer
\(L_{i+1}\), and the wqo / refutation properties propagate.

\hypertarget{theoretical-predictions}{%
\subsection{8. Theoretical Predictions}\label{theoretical-predictions}}

The framework makes the following predictions, all empirically testable:

\textbf{P1.} The sieve evicts \(\ge 30\%\) of incoming tokens at any
context length \(\ge 2048\) on natural-language workloads, under the
calibrated KSTE encoder.

\textbf{P2.} Eviction by the sieve does not increase perplexity by more
than 0.5\% on standard benchmarks (WikiText-103, C4), provided
polynomial-ring attention is retained as the scoring function.

\textbf{P3.} Ultraproduct attention with \(U\) chosen as the principal
ultrafilter at the highest-attention key reduces PPL by 1-3\% on
long-context tasks (LongBench, RULER) due to elimination of soft-mass
smearing.

\textbf{P4.} The sieve's refutation procedure (Theorem 3.1) finds
primitive-recursive counterexamples in \(\le 10^4\) operations for any
encoder-side bug, providing a debugging primitive an order of magnitude
cheaper than gradient-based attribution.

These are the falsification targets for the experimental program
described in Paper IV.

\hypertarget{relation-to-papers-i-and-ii}{%
\subsection{9. Relation to Papers I and
II}\label{relation-to-papers-i-and-ii}}

Paper I established the algebraic core: hidden state on a CM elliptic
curve over \(\mathbb{Q}(\sqrt{-163})\), Frobenius cancellation through
RMSNorm, polynomial-ring attention with KL-zero parity. Paper II
reported the system that runs that mathematics with six-figure
bit-exactness on Gemma3-1B.

This paper does \emph{not} replace Paper I or II. The sieve and
ultraproduct attention are \emph{additions} to the existing stack:

\begin{itemize}
\tightlist
\item
  \emph{Polynomial-ring attention remains the scoring function.} The
  ultraproduct construction is an alternative path that can be enabled
  selectively (\texttt{-\/-ultraproduct-attn}) but is not the default.
\item
  \emph{The Friedman sieve replaces variance-ranked top-K as the
  KV-eviction policy.} Variance-ranking is kept as a fallback
  (\texttt{-\/-legacy-eviction}).
\item
  \emph{KSTE is a side-effect of the existing Spinor block.} No
  additional memory; \textasciitilde50 LOC for tree construction.
\end{itemize}

The Frobenius framework of Paper I and the engineering scaffolding of
Paper II remain in force.

\hypertarget{open-questions}{%
\subsection{10. Open Questions}\label{open-questions}}

\begin{enumerate}
\def\labelenumi{\arabic{enumi}.}
\item
  \emph{Semantic resolution.} Is the discriminative power of \(\preceq\)
  on \(\mathcal{T}_{60,3}\) sufficient to distinguish close
  natural-language meanings? Open empirical question; Paper IV §6 is the
  experiment.
\item
  \emph{Calibration drift.} The Knight-Skeleton variance ranking is
  calibrated at warmup. Does the encoder's behaviour drift as the cache
  fills? Bounded by the residual quantizer's amax, but not yet measured.
\item
  \emph{Ultraproduct attention learnability.} Standard attention is
  differentiable; ultraproduct attention is not, directly. §11 resolves
  this by moving the learnable boundary: the choice operator \(F\) is
  part of the logical infrastructure, not a learnable parameter, and
  gradient flow lives in the \emph{generation of the class \(A\)} ---
  the structural query the model emits. The research direction
  reformulates from ``differentiable hard attention'' to
  ``differentiable class definition.''
\item
  \emph{Beyond three labels.} Is there an empirical gain from labels
  \(\{A, B, C, D\}\)? The wqo theorem holds for any finite label set,
  but the 63-byte block constraint pins us at three.
\item
  \emph{Cross-layer sieve sharing.} Can a single sieve be shared across
  heads or layers, or does each need its own? Memory implications differ
  by an order of magnitude.
\end{enumerate}

\hypertarget{the-axiomatic-layer-choice-operator-and-domain-reduction}{%
\subsection{11. The Axiomatic Layer: Choice Operator and Domain
Reduction}\label{the-axiomatic-layer-choice-operator-and-domain-reduction}}

The three operational layers of §7 --- encoder, sieve, ultraproduct
attention --- describe \emph{what} the system does. They do not yet
state \emph{why} the System-1 / System-2 boundary is well-defined. This
section provides the axiom that governs it, drawing on Friedman's
program of building consistency-strong systems from minimal logical
infrastructure.

\hypertarget{two-primitives}{%
\subsubsection{11.1 Two primitives}\label{two-primitives}}

We adopt two primitives in addition to the wqo and ultrafilter machinery
of §§3--5.

\textbf{The choice operator \(F\).} A unary function on definable
classes: \[A \ne \varnothing \;\Longrightarrow\; F(A) \in A.\] \(F\)
selects a canonical element from any nonempty class. Structurally, \(F\)
is Hilbert's \(\varepsilon\)-operator:
\(F(A) := \varepsilon x.\, x \in A\). Determinism (canonicality) is
enforced by a fixed total order \(\prec_F\) on \(\mathcal{T}_{60,3}\)
such that \(F(A) := \min_{\prec_F} A\). The implementation is specified
in Paper IV §10; the order is the packed-byte lexicographic order on the
60-byte tree representation, trivially computable and stable across
platforms.

\textbf{The Extended-Domain Reduction axiom.} Partition the universe of
cached trees into \emph{Real Objects} \(\mathrm{RO}\) (the active
context window of System 1) and an outer stratum \(\mathrm{Ext}\) (the
persistent compressed cache of System 2). For a structural predicate
\(\varphi\) and \(v \in \mathrm{RO}\),
\[\varphi(v) \;\Longrightarrow\; \varphi^*(v),\] where \(\varphi^*\) is
\(\varphi\) relativized to \(\mathrm{RO}\) --- the truth of \(\varphi\)
checked using only objects in the active window. The axiom asserts that
\emph{structural truths about the full extended domain restrict cleanly
to the active window}. Crossing the cache-to-window boundary preserves
order-relations without numerical loss.

\hypertarget{what-the-axiomatic-layer-does-for-the-architecture}{%
\subsubsection{11.2 What the axiomatic layer does for the
architecture}\label{what-the-axiomatic-layer-does-for-the-architecture}}

The layer slots in above the three operational layers without disturbing
them.

\begin{itemize}
\item
  \textbf{\(F\) is the principal-ultrafilter limit of §5.3, named
  axiomatically.} When \(U = U_p\) is the principal ultrafilter at the
  highest-ranked key position,
  \(\mathrm{UltraAttn}(Q,K,V; U_p) = V_p = F(\{V_t : t \in \mathrm{cache} \cap A\})\)
  for the structural class \(A\) derived from \(Q\). The choice operator
  and the ultraproduct limit name the same kernel; this section gives it
  the axiomatic justification the ultrafilter framing lacked.
\item
  \textbf{Extended-Domain Reduction is the attention-level analogue of
  Theorem 4 (Paper I §3.2).} Theorem 4 states: the Frobenius scale
  factor \(\pi^k\) cancels through \(QK^\top V W_O\) and vanishes at
  RMSNorm. The Reduction axiom states: structural properties of cached
  trees survive selection into the active window. Both are
  \emph{projective cancellation} theorems --- one for arithmetic scale
  through linear-algebraic operators, one for topological structure
  through the cache-to-window boundary. The two theorems together close
  the algebra: every layer of the engine has a cancellation theorem
  governing it.
\item
  \textbf{The combination resolves open question §10.3.} The
  learnability concern for ultraproduct attention was that \(F\) is not
  differentiable. Under the axiomatic framing the question dissolves:
  \(F\) is \emph{not supposed} to be a learnable parameter. It is part
  of the logical infrastructure, fixed across all inferences. The
  learnable component lives elsewhere --- in the \emph{generation of the
  structural class \(A\)}. The model learns to define classes; the
  choice operator, fixed and axiomatic, selects from them. Gradient flow
  goes to class definition, not to attention selection.
\end{itemize}

\hypertarget{inference-as-consistency-maintenance}{%
\subsubsection{11.3 Inference as consistency
maintenance}\label{inference-as-consistency-maintenance}}

Under the axiomatic layer, the inference loop is no longer ``predict the
most likely next token by minimizing cross-entropy across a continuous
distribution.'' It is closer to constraint satisfaction. Each generation
step is:

\begin{enumerate}
\def\labelenumi{\arabic{enumi}.}
\tightlist
\item
  The model defines a structural class \(A\) --- the set of trees \(T\)
  satisfying the current attention query.
\item
  The choice operator returns \(F(A)\), the \(\prec_F\)-canonical
  witness.
\item
  Extended-Domain Reduction is checked: does \(\varphi(F(A))\) imply
  \(\varphi^*(F(A))\) for the current set of invariants \(\varphi\)?
\item
  If yes, \(F(A)\) is admitted to \(\mathrm{RO}\). If no, the local
  failure is a finite, witness-bearing obstruction (§3.3, WKL₀
  refutation property) and the engine emits a correction token instead.
\end{enumerate}

For workloads where consistency is the dominant metric --- formal
verification, code generation, mathematical reasoning, structured data
extraction --- this is a strict improvement over likelihood-driven
decoding. For free-form natural-language workloads, see §11.4.

\hypertarget{fuzzy-classes-for-natural-language-workloads}{%
\subsubsection{11.4 Fuzzy classes for natural-language
workloads}\label{fuzzy-classes-for-natural-language-workloads}}

Natural language is not a logical sequence. Pure consistency-maintenance
attention may produce stilted output because the exact class \(A\) is
often too restrictive: the model wants to choose between paraphrases
that no single structural predicate distinguishes. The recovery is to
allow \(A\) to be \emph{fuzzy} --- a Hamming neighborhood of the exact
structural query, parameterized by a single radius \(r\):
\[A_r \;=\; \{T : d_{\mathrm{tree}}(T, T_Q) \le r\}.\] For \(r = 0\) the
framework reduces to hard consistency. For \(r\) large, \(A_r\)
approaches the entire cache; \(F\) then selects the \(\prec_F\)-minimum
tree subject to the soft constraint, which we may choose (by defining
\(\prec_F\) appropriately) to coincide with the top-attention key ---
recovering soft top-1 behaviour without instantiating a softmax.

The radius \(r\) is a \emph{hyperparameter, not a learnable parameter}.
It is set per-task: code generation \(r = 0\); free-form chat \(r\)
tuned for fluency; mathematical reasoning \(r = 0\) inside derivation
steps, \(r > 0\) between them. The axiomatic structure survives at every
\(r\); only the strictness of consistency varies.

\hypertarget{foundational-strength-of-the-axiomatic-layer}{%
\subsubsection{11.5 Foundational strength of the axiomatic
layer}\label{foundational-strength-of-the-axiomatic-layer}}

The choice operator with bounded comprehension is at the strength of
\(\varepsilon\)-calculus, which by Ackermann's consistency proof is
conservative over Peano Arithmetic. Extended-Domain Reduction with
structural predicates over a finite extended domain is at WKL₀-strength
by the same finitary argument as §3.3: any failure of the reduction has
a primitive-recursive witness. \textbf{The runtime stays WKL₀-strength
even with the axiomatic layer added.} The mathematics names a powerful
object; the engine executes it cheaply, and any failure remains finite
and locatable.

\hypertarget{conclusion}{%
\subsection{12. Conclusion}\label{conclusion}}

We have specified the mathematical content of the Friedman Stack. Every
layer of the stack --- encoder, sieve, ultraproduct attention, choice
operator --- rests on a primitive that admits a finite witness of
failure. The stack's expressivity reaches arbitrarily high in the
consistency-strength hierarchy (the wqo theorem alone is independent of
ATR₀); its \emph{runtime} sits at WKL₀-strength, with
primitive-recursive refutation. The asymmetry is what makes the stack
engineerable.

The companion paper (Paper IV) specifies the system, the encoder kernel,
the choice-operator implementation, the test plan, and the experimental
schedule.

\begin{center}\rule{0.5\linewidth}{0.5pt}\end{center}

\hypertarget{references}{%
\subsection{References}\label{references}}

\begin{enumerate}
\def\labelenumi{\arabic{enumi}.}
\tightlist
\item
  Kruskal, J. B., \emph{Well-Quasi-Ordering, the Tree Theorem, and
  Vázsonyi's Conjecture}, Trans. AMS 95, 1960.
\item
  Friedman, H. M., \emph{FOM archive postings} on Embedded Maximal
  Cliques (2009--2018), Foundations of Mathematics list.
\item
  Friedman, H. M., \emph{Boolean Relation Theory and Incompleteness},
  manuscript, 2014 (revised).
\item
  Friedman, H. M., \emph{Tangible Incompleteness}, manuscript series,
  2014--2024.
\item
  Harrington, L., \emph{Conservation theorem for WKL₀ over PRA},
  unpublished communication; cited in Simpson, \emph{Subsystems of
  Second Order Arithmetic}.
\item
  Simpson, S. G., \emph{Subsystems of Second Order Arithmetic},
  Cambridge UP, 2nd ed., 2009.
\item
  Gödel, K., \emph{Ontological Proof} (1970), in \emph{Collected Works
  Vol. III}, 1995.
\item
  Anderson, C. A., \emph{Some Emendations on Gödel's Ontological Proof},
  Faith and Philosophy 7, 1990.
\item
  Benzmüller, C. \& Paleo, B. W., \emph{Formalization, Mechanization and
  Automation of Gödel's Proof of God's Existence}, arXiv:1308.4526,
  2013.
\item
  Łoś, J., \emph{Quelques remarques, théorèmes et problèmes sur les
  classes définissables d'algèbres}, in \emph{Mathematical
  interpretation of formal systems}, 1955.
\item
  Shannon-Prime, \emph{Paper I --- PPT Theory} (Frobenius framework).
  2026-05-17.
\item
  Shannon-Prime, \emph{Paper II --- PPT System} (engineering on
  Gemma3-1B). 2026-05-19.
\item
  Hilbert, D. \& Bernays, P., \emph{Grundlagen der Mathematik II}, 1939
  (epsilon-operator).
\item
  Ackermann, W., \emph{Zur Widerspruchsfreiheit der Zahlentheorie},
  Math. Ann. 117, 1940 (epsilon-calculus consistency).
\item
  Friedman, H. M., \emph{FOM archive postings} on consistency proofs
  from minimal logical infrastructure (Concept Calculus and related),
  2009--2018.
\end{enumerate}

\begin{center}\rule{0.5\linewidth}{0.5pt}\end{center}

\emph{Continuation of Papers I--II. The system that executes this
paper's content is specified in Paper IV.}

\end{document}
